\documentclass[blue]{beamer}
\usepackage{beamerthemesplit}
\usepackage[spanish]{babel}
\usepackage[latin1]{inputenc}
%\usepackage{algorithm}
%\usepackage{algorithmic}
\usepackage{graphicx}
\usepackage{geometry}
\usepackage{mathrsfs}
\usepackage{amsmath}
\usepackage{pst-all}
\usepackage{amssymb}
\usepackage{url}
%\usepackage{amslatex}
\usepackage{amsmath}
\usepackage[T1]{fontenc}

\graphicspath{{./pic/}} %le dice que las im\'agenes est\'an en la carpeta de arriba

% Setup TikZ

%\usepackage{tikz}
%\usetikzlibrary{arrows}
%\tikzstyle{block}=[draw opacity=0.7,line width=1.4cm]


% Vary the color applet  (try out your own if you like)
\colorlet{structure}{red!65!black}

\beamertemplateshadingbackground{gray!50}{white}

%%%%% cuadros, sombras, redondeos, etc  %%%%%%
\mode<presentation> {
  \usetheme{Madrid}
  %\usetheme{Warsaw} %Madrid? JuanLesPins?
  % or ...

  \setbeamercovered{transparent}
  % or whatever (possibly just delete it)
}
%\logo{\includegraphics[scale=0.1]{logoUSM-DI.eps}}


\title{Proyecto \emph{V.I.Pe.R.}}
\subtitle{\texttt{Pre-Empresa {\bf Phyrex}}}
\author[Pre-Empresa {\bf Phyrex}]{%
  Rodrigo Fr\'{\i}as  \and
  Celeste Bertin      \and \\
  Patricio Carrasco   \and
  Rocio Fernandez    
}
\institute[Universidad T\'ecnica Federico Santa Maria]{Departamento de Inform\'atica \\
Universidad T\'ecnica Federico Santa Mar\'{\i}a\\
~\\
\includegraphics[scale=0.15]{logoUSM-DI.eps}
}
\date{Santiago, \today}

\begin{document}

\frame{\titlepage}

%\section[outline]{Temario}
%\frame{\tableofcontents}

%diapositiva en bordes
%\frame[plain]{}


%\section{Temario}
%\frame {
%    \frametitle{Temario}
%
%\begin{itemize}
%  \item<1-> Introducci�n.
%  \item<2-> Estado del arte.
%  \item<3-> Desarrollo.
%  \item<4-> Pruebas.
%  \item<5-> Resultados.
%  \item<6-> Conclusiones.
%  \end{itemize}
%}

\section{Explicaci\'on del Problema}
\frame {
    \frametitle{Explicaci\'on del Problema}
    \begin{block}{Razones}
      \begin{itemize}
      \item Robot Motion Planning es un problema de especificaci�n
        geom�trica.\\
      \item Por lo general se ignora aceleraci�n y velocidades angular y
        lineal.\\
      \item Idea: llegar de un punto inicial a un punto final eludiando
        obst�culos.
      \end{itemize}
    \end{block}
}

\section{Visualizaci\'on de Solucion}
\frame{
  \frametitle{Visualizaci\'on de Solucion}

}

\section{Caracter\'isticas del Segmento de Clientes}
\frame {
    \frametitle{Caracter\'isticas del Segmento de Clientes}
}

\section{Propuesta de Valor}
\frame {
    \frametitle{Propuesta de Valor}
}

\section{Fundamentaci\'on de la Innovaci\'on}
\frame {
    \frametitle{Fundamentaci\'on de la Innovaci\'on}
}
\end{document}
